\scanpage{29}
% slanted
% loc: upper left slant
\margin{$(\unsure{\ast\ast|\ast})$ \unsure{exprime|expriment|exprimons|exigeons} que si $a\ne b$, alors $\pi_a\ne\pi_b$}
% loc: 1–4
De plus, si $a_1,a_2\in A$, $a_1\ne a_2$, on voit de suite que les décompositions $\pi_{a_1},\pi_{a_2}$ sont bien parallèles, à cela près qu'il faut exiger \deleted{\unsure{que si $s$ est un sommet d'or}}%
% loc: line 4, below deletion and continued in left margin
\inserted{$(\unsure{\ast\ast})$ si $s\in S$ est d'ordre $2$, alors un des trois \unsure{ens|ens.} $U_s,E_a,E_b$ est $\ne\varnothing$, où $a,b$ sont les deux arêtes \unsure{inc.|inci.} à $s$}

% loc: 5–8
Une \unsure{\deleted{façon}|façon} de visualiser la construction est via un arbre topologique $\unsure{\mathcal{X}|\mathscr{X}|X}\deleted{\unsure{A|\Gamma}}=(\overset{\text{\inserted{\unsure{esp. top.}}}}{\underline{\Gamma}},S)$ (tel que $\Gamma$ en soit une description combinatoire). On fixe un pt sur
% loc: 9–12
chaque \unsure{arête|arète}, de façon à \unsure{donner|trouver} un graphe topologique (subdivision barycentrique) $\mathcal{X}^{\ast}=(\underline{\Gamma}^{\unsure{\ast}},\underbrace{S\amalg A'}_{S^{\ast}})$, où $A'\unsure{\subset|\in}\underline{\Gamma}$ est en \unsure{corr|corr.} \deleted{\unsure{b|bi}}biunivoque donnée avec \deleted{\unsure{tout}} $A$.

% loc: 13–17, bracketed NB
\deleted{\unsure{pt de}} [\uline{NB} $A'$ est formé de \deleted{pts} sommets d'ordre $2$, et \inserted{pour} \unsure{tte|toute} arête de $\mathcal{X}^{\ast}$, une exactement des deux extrémités appartient à $A'$, l'autre à $S$ -- voici une description complète des propriétés particulières \unsure{aux|des} $\underline{\Gamma}^{\ast}$\unsure{..|...}]

% loc: 18–21
Alors \deleted{les}\unsure{le} point\deleted{s} \inserted{$a'$} de $A'$ coupe\deleted{\unsure{nt}} $\underline{\Gamma}^{\deleted{\unsure{\ast}}}$ en deux composantes connexes, dont les adhérences définissent une décomposition binaire de $\underline{\Gamma}$, d'\unsure{``|''}équateur $\{a'\}$. Deux
% loc: 22–23
telles décompositions \deleted{\unsure{bi}} de $\underline{\Gamma}$, pour $a'\ne\unsure{a''|a'}$, définissent des décompositions binaires \unsure{strict$^t$|stric$^t$} non parallèles.
% archive sheet number: 28
